% Part: many-valued-logic
% Chapter: sequent-calculus
% Section: introduction

\documentclass[../../../include/open-logic-section]{subfiles}

\begin{document}

\olfileid{mvl}{seq}{int}

\olsection{పరిచయం}

సాంప్రదాయిక తర్కంలోని సీక్వెంట్ కలనం సమర్థమైన, సరళమైన
\tetoken{వ్యుత్పత్తి}{!!{derivation}} వ్యవస్థ. ఒక బహు-విలువల తర్కం
పరిమిత సంఖ్యలో సత్యమూల్యాలు గల మాత్రికచే నిర్వచించబడితే,
అంటే $V$ పరిమిత సమితి అయితే, దానికి సీక్వెంట్ కలనం ఇవ్వవచ్చు.
ఎలా ఇవ్వాలనే ఆలోచన సాంప్రదాయిక సందర్భంలో సీక్వెంట్ల
అర్థాలను, నిగమన నియమాల రూపాన్ని పరిశీలించడం నుంచి వస్తుంది.

ఒక సీక్వెంట్‌ను ఈ విధంగా గుర్తుచేసుకుందాం:
\begin{align*}
    !A_1, \dots, !A_m & \Sequent !B_1, \dots, !B_n
\intertext{దీనిని కింది \tetoken{సూత్రంగా}{!!{formula}} అర్థం చేసుకోవచ్చు:}
    (!A_1 \land \cdots \land !A_m) & \lif (!B_1 \lor \cdots \lor
!B_n)
\end{align*}
\sourcecorrection{OLTEMVLSEQ-001}{ఆంగ్ల మూలంలోని
సీక్వెంట్ ఎడమవైపు చివరి సూచిక $n$గా ఉంది; వెంటనే
ఇచ్చిన దానికి సరిపడే సూత్రంలో అది $m$.
రెండు వైపుల సంఖ్యలను వేరుగా చూపడానికి
సీక్వెంట్‌లోని చివరి ఎడమ సూచికను $m$గా సరిచేశాం.}
మరో మాటలో, \tetoken{సత్యమూల్య కేటాయింపు}{!!^a{valuation}}~$\pAssign{v}$
సీక్వెంట్ $\Gamma \Sequent \Delta$ను \emph{సంతృప్తిపరుస్తుంది},
ఏదో $!A \in \Gamma$కు $\pValue{v}(!A) = \False$
లేదా ఏదో $!A \in \Delta$కు $\pValue{v}(!A) = \True$
అయితే, అప్పుడు మాత్రమే. ఈ అర్థం ప్రకారం ప్రారంభ
సీక్వెంట్లు $!A \Sequent !A$ ఎప్పుడూ సంతృప్తిపడతాయి:
$\pValue v(!A) = \True$ లేదా $\pValue v(!A) =
\False$ తప్పనిసరిగా ఉంటుంది.
\sourcecorrection{OLTEMVLSEQ-002}{చివరి మూల సమీకరణలో
$\pValue$కు కేటాయింపు $v$ ఆర్గ్యుమెంట్ తప్పిపోయింది.
అదే పేరాలోని ముందరి రెండు మూల్యాంకనాల
మాదిరిగా దాన్ని చేర్చాం; సత్య షరతు మారలేదు.}

$\Gamma$, $\Delta$ అనే పక్క సూత్రాలను చూపకుండా,
$\Log{LK}$లోని సోపాధికానికి నిగమన నియమాలు ఇవి:

\begin{defish}
    \Axiom$ \fCenter !A$
    \Axiom$ !B \fCenter $
    \RightLabel{\LeftR{\lif}}
    \BinaryInf$ !A \lif !B \fCenter $
    \DisplayProof
    \hfill
    \Axiom$ !A \fCenter !B$
    \RightLabel{\RightR{\lif}}
    \UnaryInf$ \fCenter !A \lif !B $
    \DisplayProof
\end{defish}

సీక్వెంట్‌కు పైన చెప్పిన అర్థపర వివరణను
వర్తింపజేస్తే, $\LeftR{\lif}$ నియమం ఇలా
చెబుతుందని చదవవచ్చు: $\pValue v(!A) = \True$,
$\pValue v(!B) = \False$ అయితే
$\pValue v(!A \lif !B) = \False$.
అదేవిధంగా, $\RightR{\lif}$ నియమం
$\pValue v(!A) = \False$ లేదా
$\pValue v(!B) = \True$ అయితే
$\pValue v(!A \lif !B) = \True$ అని
చెబుతుంది. నిజానికి ఈ రెండు షరతులూ
ద్విదిశ షరతులే. $\LeftR{\land}$,
$\RightR{\lor}$ నియమాల ప్రామాణిక
రూపాల్లో వాటికి సరిపోలే షరతులు
ద్విదిశ షరతులు కావు. అయితే,
అలా అయ్యే ప్రత్యామ్నాయ నియమరూపాలు ఇవి:

\begin{defish}
    \Axiom$!A, !B, \Gamma \fCenter \Delta$
    \RightLabel{\LeftR{\land}}
    \UnaryInf$!A \land !B, \Gamma \fCenter \Delta$
    \DisplayProof
    \hfill
    \Axiom$ \Gamma \fCenter \Delta, !A, !B$
    \RightLabel{\RightR{\lor}}
    \UnaryInf$ \Gamma \fCenter \Delta, !A \lor !B$
    \DisplayProof
\end{defish}

ఈ ప్రాథమిక ఆలోచనను $n$-విలువల తర్కానికి
వర్తింపజేస్తే, రెండు స్థానాలకు బదులు $n$
స్థానాలున్న సీక్వెంట్ కలనం వస్తుంది;
ప్రతి సత్యమూల్యానికి ఒక స్థానం ఉంటుంది.
$V = \{\False, \Undef, \True\}$ గల
మూడు-విలువల తర్కంలో సీక్వెంట్
$\Gamma \mid \Pi \mid \Delta$ రూపంలో ఉంటుంది.
\tetoken{సత్యమూల్య కేటాయింపు}{!!a{valuation}}~$\pAssign v$
దాన్ని సంతృప్తిపరుస్తుంది, ఏదో
$!A \in \Gamma$కు $\pValue{v}(!A) = \False$
లేదా ఏదో $!A \in \Delta$కు
$\pValue{v}(!A) = \True$ లేదా ఏదో
$!A \in \Pi$కు $\pValue{v}(!A) = \Undef$
అయితే, అప్పుడు మాత్రమే. అందువల్ల
ప్రారంభ సీక్వెంట్లు $!A \mid !A \mid !A$
ఎప్పుడూ సంతృప్తిపడతాయి.

\end{document}
